\documentclass[10pt,a4paper]{amsart}
\usepackage[margin=19mm]{geometry}
\usepackage{amsmath,amssymb,mathtools}
\usepackage[scr=rsfs,cal=euler]{mathalfa}
\usepackage{xfrac}
\usepackage[unicode,hidelinks,bookmarks=false]{hyperref}
\newcommand{\ie}{i.e.\ }
\newcommand{\cf}{cf.\ }
\newcommand{\resp}{respectively\ }
\newcommand{\Subsection}{\subsection}
\providecommand{\nobreakdash}{\nobreak\hbox{-}}
\setlength{\emergencystretch}{2em}
\multlinegap0pt
\usepackage{mathtools}
\usepackage{amssymb}
\usepackage{tikz}
\newtheorem{lemma}{Lemma}
\title{Real-rootedness of the $\tau$-polynomial under graph joins}
\author{Mingyang Kang, Zhixin Liu, Sophie C.C. Sun, Philip B. Zhang}
\date{Source: arXiv:2609.07457v1 (CC BY 4.0)}
\begin{document}
\maketitle

\noindent\emph{Source and licence.} Real-rootedness of the $\tau$-polynomial under graph joins. Version arXiv:2609.07457v1 (CC BY 4.0), distributed by arXiv under the Creative Commons Attribution 4.0 International licence, http://creativecommons.org/licenses/by/4.0/, as recorded in the arXiv licence metadata for this exact version and shown on the abstract page https://arxiv.org/abs/2609.07457v1. Full licence notice: \texttt{licenses/arxiv-2609.07457-CC-BY-4.0.txt}.\par\smallskip

\noindent\textit{Editorial scope.} Counted unit: the article's lemma lem:coefficient-substitution (source0.tex lines 523--543). The article's complete original proof of it is delivered with it (source0.tex lines 545--613, one \texttt{proof} environment of 69 lines). No mathematics is written, repaired or re-derived: the statement and the proof are the article's own lines, in the article's own notation, and no retained line is substituted or re-flowed.  No further source-local context is needed: every cross-reference of every retained line resolves inside the two delivered spans.  Nothing was omitted from any retained span, so the package carries no omission record.  No retained line contains a citation, so the wrapper carries no bibliography and no bibliography entry.  No retained line contains a figure environment, an image-inclusion command or drawing code, so no image is delivered and the package declares no asset dependency.  Editorial formatting only, in addition: an AMS \texttt{amsart} preamble that restates the article's own declarations and macros used by the retained lines, namely the environments \texttt{lemma} on separate counters, together with the standard packages those lines need.  The retained environments print in this document as Lemma 1 in order of appearance, whereas the article prints the same items with the numbering of its own full document.  The complete native source of the article as distributed in the arXiv e-print for this exact version is bundled unchanged as \texttt{sources/209629/source0.tex}.
\begin{lemma}
\label{lem:coefficient-substitution}
Let $ D=1-uv$, and 
\begin{equation}\label{nj-9-mk9}
    \xi(u,v)=\frac{u(1+v)}{D},\qquad
 \eta(u,v)=v(1+\xi(u,v))=\frac{v(1+u)}{D},
\end{equation}
with Jacobian determinant
\begin{equation}\label{bn-sn-09}
 J(u,v)=\det\frac{\partial(\xi,\eta)}{\partial(u,v)}
       =\frac{(1+u)(1+v)}{D^3}.    
\end{equation}
Then, for every formal power series $F(\xi,\eta)$ and all $r,s\geq0$,
\begin{equation}\label{eq:coefficient-substitution}
 [u^rv^s]\,J(u,v)F(\xi,\eta)
 = [\xi^r\eta^s]
   (1+\xi)^{s+1}(1+\eta)^{r+1}F(\xi,\eta),
\end{equation}
where, on the right-hand side, $\xi$ and $\eta$ are treated as independent
  variables.
\end{lemma}

\begin{proof}
We first need  a one-variable identity.  Suppose that
\[
 z=\frac{ct}{1-dt}, 
\]
where $c$ is invertible.  Let $H(z)=\sum_{n\geq0}h_nz^n$ be a formal power series.  Then, for
$r\geq0$,
\begin{equation}\label{eq:mobius-coefficient}
 [t^r]\frac{\partial z}{\partial t}\cdot H(z)
 =[z^r](c+dz)^{r+1}H(z).
\end{equation}
This can be justified as follows.  Since
\[
 \frac{\partial z}{\partial t}=\frac{c}{(1-dt)^2},
\]
the contribution of each $h_n$ ($0\leq n\leq r$) to the left-hand side of \eqref{eq:mobius-coefficient} is
\[
 h_n[t^{r-n}]\frac{c^{n+1}}{(1-dt)^{n+2}}
 =h_n\binom{r+1}{r-n}c^{n+1}d^{r-n},
\]
which exactly agrees with the  contribution of $h_n$ to the right-hand side  of \eqref{eq:mobius-coefficient}. 
This verifies \eqref{eq:mobius-coefficient}. 

We now apply \eqref{eq:mobius-coefficient} to prove \eqref{eq:coefficient-substitution}.
Set  $c=1+v$, $d=v$, and 
\begin{equation}\label{x-12}
     z=\frac{cu}{1-du}=\xi.
\end{equation}
Note also that 
\[
v(1+z)=\eta.
\] 
Let 
\begin{align}
     H(z)&=(1+z)F\bigl(z,v(1+z)\bigr)\nonumber\\
     &=(1+\xi)F(\xi, \eta).\label{x-09}
\end{align}
Applying 
\eqref{eq:mobius-coefficient} with $t$ replaced by the variable $u$,
\begin{equation}\label{eq:mobius-coe}
 [u^r]\frac{\partial z}{\partial u}\cdot H(z)
 =[z^r](c+dz)^{r+1}H(z).
\end{equation}

By \eqref{x-12} and \eqref{x-09}, one may compute that 
\begin{equation*}
  \frac{\partial z}{\partial u}\cdot H(z)=J(u,v) F(\xi,\eta).  
\end{equation*}
On the other hand, noting that 
  $c+dz=1+\eta$, we have
\begin{equation*}
    (c+dz)^{r+1}H(z)=(1+\eta)^{r+1}(1+\xi)F(\xi,\eta),
\end{equation*}  
{where $\eta=v(1+\xi)$}. 
So \eqref{eq:mobius-coe} can be rewritten as 
\begin{equation}\label{eq:first-coefficient-step}
 [u^r]J(u,v)F(\xi,\eta)
 =[\xi^r](1+\eta)^{r+1}(1+\xi)F(\xi,\eta),
\end{equation}
where, on the right-hand side,  $\eta=v(1+\xi)$ is regarded as a function in~$\xi$. 


Finally, for every series $G(\eta)$, note that 
\[
 [v^s]G\bigl(v(1+\xi)\bigr)
 =(1+\xi)^s[\eta^s]G(\eta).
\]
Therefore, taking $[v^s]$ in \eqref{eq:first-coefficient-step}   yields \eqref{eq:coefficient-substitution}.
\end{proof}

\end{document}
